%%
% 第六章 结论与展望
% 短且克制。四点创新 + 局限性如实声明 + 未来工作仅限方法学与数据层面。
%%

\chapter{结论与展望}
\label{cha:conclusion}

\section{主要结论与创新贡献}
\label{sec:contributions}

本论文以 \textit{Nature}~2025 发表的酶底物特异性预测 SOTA 模型 EZSpecificity\cite{EZSpec2025} 为参照框架，以其训练数据集 ESIBank 中覆盖显著稀缺（$389/25{,}225 \approx 1.5\%$）但应用价值极高的细胞色素 P450 超家族为试金石，系统检验了通用 SOTA 模型在稀缺家族上的 对未见酶的泛化能力，并在领域专属数据上开展重训与架构改进。综合全文工作，本研究的主要结论与创新贡献可归纳为以下四点。

\subsection{构建了目前规模最大、标注最详细的 P450 酶--底物特异性公开基准数据集}
\label{subsec:contrib-dataset}

现有 P450 数据库存在普遍局限：RCSB PDB\cite{Burley2025,Berman2000} 中共结晶配体可能是底物、产物、抑制剂或辅因子，未经区分直接使用会引入大量假阳性；ESIBank\cite{EZSpec2025} 基于 EC 号粗筛导致 FAD 单加氧酶等非 P450 混入；P450Rdb 的 132{,}000 条原始记录实为两表拼接，仅 3{,}753 条为实验验证条目；PlantP450DB\cite{Nelson2004} 与 PCPD\cite{Wang2021PCPD} 大量条目仅含酶序列与笼统反应描述而无标准化底物 SMILES。针对上述局限，本研究通过文献检索（PubMed + CrossRef）、PDB 结构分析（Fe--配体距离、配位关系、修饰特征）与多轮人工审核（辅以多智能体协作流程对每条酶--配体记录证据复核）严格区分底物、产物、抑制剂、辅因子与溶剂类型，整合得到 $1{,}622$ 个唯一酶、$2{,}125$ 个唯一底物、$4{,}751$ 条跨源去重正样本对、$47{,}510$ 条可用训练评估样本；相比 ESIBank 原始训练集在 P450 上的覆盖规模扩大约 4 倍，并显著扩展了植物 P450 的数据覆盖（PlantP450DB 与 PCPD 合计贡献 $2{,}188$ 条正样本与 $1{,}396$ 个唯一酶）。该基准数据集与对应的证据复核流程为后续社群验证、扩展与湿实验选样提供了可复用的基础设施。

\subsection{建立了一套泄漏控制的外部公平对比评估协议}
\label{subsec:contrib-protocol}

针对 SOTA 模型外部评估中训练集--测试集酶重合引发的记忆效应问题，本研究以 InterPro\cite{Blum2025} 功能域对 ESIBank 原始 $25{,}225$ 条酶严格识别出 $389$ 个真正的 P450，并据此从本研究测试集中剔除与之重合的 $3{,}036$ 条样本，得到 $7{,}963$ 条过滤后外部测试集。该协议将通用模型在 P450 上的评估条件收敛为"完全未见过的 P450 酶"，是对 SOTA 通用性声明在稀缺家族上进行严格检验的必要前提。

\subsection{实证发现 SOTA 通用模型在数据缺乏家族上存在未见酶的泛化瓶颈、领域专属重训可显著突破该瓶颈}
\label{subsec:contrib-finding}

在过滤后外部测试集（$n=7{,}963$）上：EZSpecificity 公开 checkpoint 仅取得 Test AUC $= 0.5586$、Test AUPR $= 0.1004$，AUPR 接近样本正负比下的随机基线；本研究在相同架构与领域专属数据上从头训练得到的基线模型 M0 达到 Test AUC $= 0.9205$、Test AUPR $= 0.6403$，相对参考模型的差距为 $\Delta\mathrm{AUC} = +0.3619$ 与 $\Delta\mathrm{AUPR} = +0.5399$。这一经验观察为"面向训练集稀缺家族的领域专属数据微调是必要的"这一方法学命题提供了量化支撑。

\subsection{提出并实现了面向方向信息感知的双图架构改进，并在单 fold 评估下展现了初步可行性}
\label{subsec:contrib-arch}

在原架构的原子级 EGNN 通路之外，本研究新增了一条残基级几何向量感知（GVP）通路，通过残基骨架二面角标量与 $\mathrm{C}\alpha \to \mathrm{N/C/C}\beta$ 向量联合编码捕获 SE(3)-等变的方向信息，并以零初始化 $+$ 延迟解冻的融合设计保证训练起步时对基线的严格等价（第 0 步数值最大绝对偏差 $4 \times 10^{-8}$）。在单 random fold 评估下，于未过滤测试集（$n=10{,}999$）上：双图架构 M1 相对原架构 M0 观察到 Test AUC $0.9302 \to 0.9310$ 的微幅提升（$\Delta = +0.0008$），同时 Test AUPR 由 $0.6749$ 回落至 $0.6474$（$\Delta = -0.0275$），两项指标并列对称报告。该观察作为方向信息感知设计的初步可行性探索，不作为性能超越或稳健性结论。

\subsection{研究意义与作用}
\label{subsec:contrib-impact}

本研究为 P450 底物特异性预测提供了可供社群验证与扩展的领域专属基准与泄漏控制评估协议，示范了在进行家族专属深度学习模型评估时泄漏控制的重要性；所提双图架构改进方案为后续在 P450 领域探索几何感知模型设计提供了可行的技术路线。研究结果对药学中 P450 代谢筛选工具的可信度评估、以及合成生物学中植物 P450 选择的辅助决策具有方法学参考意义。

\section{局限性}
\label{sec:limitations-final}

本研究的结论均严格限定于以下若干可审计的实验范围，相应局限在此集中声明。

\subsection{仅在单一随机切分上完成训练与评估}
\label{subsec:limit-fold}

本研究全部训练与测试结果均基于 random split 第一折，未开展 $k$ 折交叉验证、多随机种子重复或置信区间估计。因此第四章 $+0.3619$ AUC 差值与第五章 $+0.0008$ AUC 差值的稳健性范围未被系统量化。

\subsection{结构增强的观察不足以支撑稳健性断言}
\label{subsec:limit-structure}

双图架构 M1 的 AUC 微幅提升（$+0.0008$）处于单 fold 评估的噪声边缘，且 Test AUPR 同步出现 $-0.0275$ 的回落。本论文如实并列报告两项变化，将其定位为架构工程层面的初步可行性探索，不作为性能改进或机制结论；稳健性需待多 fold $\times$ 多 seed $\times$ 超参充分调优 $\times$ 显著性检验的系统评估。

\subsection{负样本基于随机配对假设}
\label{subsec:limit-neg}

本研究遵循 EZSpecificity 的双向随机负样本构造策略，假设在本研究样本空间中随机配对的酶--底物对以高概率为非催化关系。P450 家族内酶之间的序列与结构相似性较高，随机负样本中可能混入未被记录的真实正配对，引入假阴性；同家族难负样本的受控构造与评估未在本论文中系统开展。

\subsection{外部泄漏过滤仅限于酶层面}
\label{subsec:limit-leak}

本研究的过滤协议剔除了与 ESIBank 原始训练集中 $389$ 个 P450 酶重合的测试样本，但未进一步过滤底物层面与 ESIBank 训练集重合的部分；若某些底物作为 ESIBank 训练集中其他家族酶的底物出现过，其表征分布对参考模型可能仍具有一定熟悉性。

\subsection{对接姿态噪声与结构特征的不确定性}
\label{subsec:limit-dock}

对接姿态与真实结合构象之间存在不可避免的误差，结构通道的输入质量因此存在天然上限。AlphaFill 回填血红素的预测结构与晶体结构在活性位点细节上的一致性未做逐条评估。

\section{未来改进方向}
\label{sec:future-work}

本研究为 P450 底物特异性预测提供了数据、评估与架构三个层面的基础材料，后续工作可沿以下方法学与数据层面的方向展开。

\subsection{更严格的切分与稳健性评估}
\label{subsec:future-split}

在现有数据集上开展多随机切分 $+$ $k$ 折交叉验证 $+$ 多种子重复以建立稳健性置信区间；按酶的序列相似性聚类设计阶梯化测试集（例如限定测试与训练的序列一致度 $<40\%$），使评估难度梯度化并暴露模型在不同同源距离下的泛化曲线。

\subsection{困难负样本与评估协议的扩展}
\label{subsec:future-neg}

针对"同家族难负样本"单独构造可控测试集，系统度量模型在同家族内部区分不同酶--底物偏好的能力；将外部泄漏过滤从酶层面扩展到底物层面以进一步收紧公平对比条件。

\subsection{底物类别作为辅助监督信号}
\label{subsec:future-class}

本研究在数据清洗过程中已将 $2{,}125$ 个底物按 NPClassifier\cite{NPClassifier2021} 体系分为萜类、氨基酸、脂肪酸、生物碱、甾体、苯丙素、聚酮与其他共 $8$ 大类；未来可将底物类别作为辅助任务监督主模型的多任务训练，检验类别先验对 P450 特异性预测的帮助。

\subsection{双图架构的稳健性验证与拓展}
\label{subsec:future-arch}

在多 fold 稳健性评估框架下重做双图架构 M1 与原架构 M0 的对比，以确认本论文观察到的 $+0.0008$ AUC 提升与 $-0.0275$ AUPR 回落分别是否具有跨 fold 的一致性；在此基础上探索融合策略（深融合 vs 浅旁路权重）与 GVP 层数等超参的系统扫描。

\subsection{从二分类到区域选择性与反应级预测}
\label{subsec:future-reaction}

P450 的实际价值不仅在于判断酶--底物配对的是非，更在于预测催化后的产物结构与区域选择性；本研究数据集中的 PCPD 反应图片（$857$ 张）与 P450Rdb 反应 SMILES 可作为向反应级预测扩展的基础数据。

\subsection{可解释性与湿实验验证}
\label{subsec:future-validate}

在模型预测基础上做注意力热图与口袋--底物原子级贡献分析以辅助药学家解释；选择若干高置信度预测配对开展体外酶活实验或 KCAT/Km 测定以验证模型可靠性，并反馈用于模型迭代。
